\documentclass[10pt,a4paper]{amsart}
\usepackage[margin=19mm]{geometry}
\usepackage{amsmath,amssymb,mathtools}
\usepackage[scr=rsfs,cal=euler]{mathalfa}
\usepackage{xfrac}
\usepackage[unicode,hidelinks,bookmarks=false]{hyperref}
\newcommand{\ie}{i.e.\ }
\newcommand{\cf}{cf.\ }
\newcommand{\resp}{respectively\ }
\newcommand{\Subsection}{\subsection}
\providecommand{\nobreakdash}{\nobreak\hbox{-}}
\setlength{\emergencystretch}{2em}
\multlinegap0pt
\usepackage{amssymb}
\usepackage{mathtools}
\usepackage{xcolor}
\usepackage{dsfont}
\newtheorem{lemma}{Lemma}
\newcommand{\De} {\Delta}
\newcommand{\La} {\Lambda}
\newcommand{\Nh}{\mathcal{N}_\la}
\newcommand{\Q}{\textit{Q}}
\newcommand{\R}{\mathbb{R}}
\newcommand{\Th} {\Theta}
\newcommand{\X}{\mathcal{X}_{0}}
\newcommand{\ds} {\displaystyle}
\newcommand{\ga} {\gamma}
\newcommand{\grad}{\nabla}
\newcommand{\la} {\lambda}
\newcommand{\lb}{\left(}
\newcommand{\na} {\nabla}
\newcommand{\om} {\Omega}
\newcommand{\ra} {\rightarrow}
\newcommand{\rb}{\right)}
\newcommand{\rnn}{\mathbb{R}^{N}}
\newcommand{\ve}{\varepsilon}
\title{Multiplicity results for mixed local-nonlocal variable exponent problem involving singular and superlinear term}
\author{Shammi Malhotra, Ambesh Kumar Pandey, K. Sreenadh}
\date{Source: arXiv:2506.19793v2 (CC BY 4.0)}
\begin{document}
\maketitle

\noindent\emph{Source and licence.} Multiplicity results for mixed local-nonlocal variable exponent problem involving singular and superlinear term. Version arXiv:2506.19793v2 (CC BY 4.0), distributed by arXiv under the Creative Commons Attribution 4.0 International licence, http://creativecommons.org/licenses/by/4.0/, as recorded in the arXiv licence metadata for this exact version and shown on the abstract page https://arxiv.org/abs/2506.19793v2. Full licence notice: \texttt{licenses/arxiv-2506.19793-CC-BY-4.0.txt}.\par\smallskip

\noindent\textit{Editorial scope.} Counted unit: the article's lemma lm:existence\_result (source0.tex lines 3677--3681). The article's complete original proof of it is delivered with it (source0.tex lines 3684--3948, one \texttt{proof} environment of 265 lines). No mathematics is written, repaired or re-derived: the statement and the proof are the article's own lines, in the article's own notation, and no retained line is substituted or re-flowed.  Retained uncounted as the source-local context the proof needs: lines 140--149; lines 767--776; lines 2871--2884; lines 3234--3243; lines 3522--3532.  Nothing was omitted from any retained span, so the package carries no omission record.  No retained line contains a citation, so the wrapper carries no bibliography and no bibliography entry.  No retained line contains a figure environment, an image-inclusion command or drawing code, so no image is delivered and the package declares no asset dependency.  Editorial formatting only, in addition: an AMS \texttt{amsart} preamble that restates the article's own declarations and macros used by the retained lines, namely the environments \texttt{lemma} on separate counters, together with the standard packages those lines need.  The retained environments print in this document as Lemma 1 in order of appearance, whereas the article prints the same items with the numbering of its own full document.  The complete native source of the article as distributed in the arXiv e-print for this exact version is bundled unchanged as \texttt{sources/180449/source0.tex}.
	\begin{equation}\label{eq1}\tag{$\mathcal{Q}_{\lambda}$}
   \left\{
	\begin{aligned}
		-\De_{p(x)}u+\lb -\De\rb^s_{p(x)}u 
&=\lambda a(x)u^{-\gamma(x)}+b(x)u^{r(x)-1},\,
u>0 \ \text{ in } \om,\\
u &= 0 \qquad \text{ in } \rnn\backslash\om,
	\end{aligned}
    \right.
\end{equation}

\begin{equation}\label{eq2}\tag{$\mathcal{Q}^+_{\lambda}$}
   \left\{
	\begin{aligned}
		-\De_{p(x)}u+\lb -\De \rb^s_{p(x)}u 
&=\lambda a(x)u_+^{-\gamma(x)}+b(x)u_+^{r(x)-1},\,
 u>0 \  \text{ in } \om,\\
u &= 0 \qquad \text{ in } \rnn\backslash\om.
	\end{aligned}
    \right.
\end{equation}

\begin{lemma}\label{lm9}
    Let $u\in \mathcal{N}_\la^\pm,$
then there exists $\varepsilon>0$
and a continuous function 
$f:B_\varepsilon(0)\ra \R^+$
such that
$$
    f(u)>0, f(0)=1 \text{ and } 
    f(v)(u+v)\in \mathcal{N}_\la^\pm 
    \text{ for all } v\in B_\varepsilon(0),
$$
    where $B_\varepsilon(0):=
    \{v\in \X:\|v\|_{\X}<\varepsilon\}.$
\end{lemma}

\begin{lemma}\label{lm:min_neg}
    For $0<\la<\La,$
then $J_\la$
has a minimizer $z_0\in \Nh^-$
and satisfies the following
$$
    J_\la(z_0)=
    \inf_{u\in \mathcal{N}_\la^-}J_\la(u)=\Th_\la^->0.
$$
\end{lemma}

\begin{align}\label{eq16}
&
\int_\om|\grad u|^{p(x)-2}\na u\na\phi \notag \\
& +
 \iint_{\mathbb{R}^{2N}} 
\frac{|u(x)-u(y)|^{p(x,y)-2}(u(x)-u(y))
(\phi(x)-\phi(y))}{|x-y|^{N+sp(x,y)}}\nonumber\\
    &-\la
\int_\om a(x)u_+^{-\ga(x)}\phi -
\int_\om b(x)u_+^{r(x)-1}\phi\ge0.
    \end{align}

\begin{lemma}\label{lm:existence_result}
    The functions $u$
and $v$
are positive weak solutions of \eqref{eq1}.
\end{lemma}

\begin{proof}[\bf Proof.]
    For any $\phi\in \X$
and $\ve>0,$
set $\Phi(x)=(u+\ve\phi)_+(x).$
Let $\om=\om_1\times\om_2$
where
$$
    \om_1:=\{x\in\om:u+\ve\phi>0\}
    \text{ and } 
    \om_2:=\{x\in\om:u+\ve\phi\le0\}.
    $$
    We can see that $\Phi|_{\om_1}=u+\ve\phi$
and $\Phi|_{\om_2}=0.$
Next, we decompose $\Q$
in the following way
\begin{align*}
\Q:= (\om_1\times\om^c) 
& \cup(\om_2\times\om^c)
\cup(\om^c\times\om_1)  
  \cup (\om^c\times\om_2) \\ &
\cup(\om_2\times\om_1)\cup
 (\om_1\times\om_2)   \cup(\om_1\times\om_1)\cup
  (\om_2\times\om_2).
\end{align*}
For convenience, 
we denote 
$$
d \nu := \frac{dx \, dy}{|x-y|^{N+sp(x,y)}},
$$
and
$$
W(x,y):=U(x,y)[(u+\ve\phi)_-(x)-(u+\ve\phi)_-(y)] ,
$$
where $U(x,y):=|u(x)-u(y)|^{p(x,y)-2}(u(x)-u(y))$.
% and $d \nu = \frac{dx \, dy}{|x-y|^{N+sp(x,y)}}$
% $$
% U(x,y):=|u(x)-u(y)|^{p(x,y)-2}(u(x)-u(y)) 
% \  \text{ and $K(x,y):=
% \frac{1}{|x-y|^{N+sp(x,y)}}.$}
% $$

Thus, we have
\begin{enumerate}
\item[(i)]$\ds
\int_{\om_1\times\om^c}W(x,y)\,d\nu=
\int_{\om^c\times\om_1}W(x,y)\,d\nu=0.$

\item[(ii)]$\ds 
\int_{\om_2\times\om^c}W(x,y)\,d\nu=-
\int_{\om_2\times\om^c}|u(x)|^{p(x,y)-2}u(x)
(u+\ve\phi)(x) \,d\nu.$

\item[(iii)]$\ds 
\int_{\om^c\times\om_2}W(x,y)\,d\nu=-
\int_{\om_2\times\om^c}|u(x)|^{p(x,y)-2}u(x)
(u+\ve\phi)(x) \,d\nu.$

\item[(iv)]$\ds 
\int_{\om_2\times\om_1}W(x,y)\,d\nu=-
\int_{\om_2\times\om_1}U(x,y)(u+\ve\phi)
(x) \,d\nu.$

\item[(v)]$\ds 
\int_{\om_1\times\om_2}W(x,y)\,d\nu=-
\int_{\om_2\times\om_1}U(x,y)
(u+\ve\phi)(x) \,d\nu.$

\item[(vi)]$\ds 
\int_{\om_1\times\om_1}W(x,y)\,d\nu=0.$

\item[(vii)]$\ds 
\int_{\om_2\times\om_2}W(x,y)\,d\nu=-
\int_{\om_2\times\om_2}U(x,y)((u+\ve\phi)(x)
-(u+\ve\phi)(y)) \,d\nu.$
\end{enumerate}
Substituting $\Phi$
in \eqref{eq16}, we get
\begin{align*}
    0\le&
\int_\om|\grad u|^{p(x)-2}\na u\na\Phi+
\int_{\Q} 
\frac{U(x,y)(\Phi(x)
-\Phi(y))}{|x-y|^{N+sp(x,y)}}\nonumber -\la
\int_\om a(x)u_+^{-\ga(x)}\Phi \\
&-\int_\om b(x)u_+^{r(x)-1}\Phi\\
    =&
\int_\om|\grad u|^{p(x)-2}\na u\na(u+\ve\phi)+
\int_{\Q} 
\frac{U(x,y)((u+\ve\phi)(x)-
(u+\ve\phi)(y))}{|x-y|^{N+sp(x,y)}}\\
    &-\la
\int_\om a(x)u_+^{-\ga(x)}(u+\ve\phi) -
\int_\om b(x)u_+^{r(x)-1}(u+\ve\phi) \\
 &+
\int_\om|\grad u|^{p(x)-2}\na u\na(u+\ve\phi)_- 
+
\int_{\Q} 
\frac{W(x,y)}
{|x-y|^{N+sp(x,y)}} \\
&   -\la
\int_\om a(x)u_+^{-\ga(x)}(u+\ve\phi)_- -
\int_\om b(x)u_+^{r(x)-1}(u+\ve\phi)_-\\
    =&\ve \left\{
\int_\om|\grad u|^{p(x)-2}\na u\na\phi+
\int_{\Q} 
\frac{U(x,y)(\phi(x)-\phi(y))}{|x-y|^{N+sp(x,y)}}-\la
\int_\om a(x)u_+^{-\ga(x)}\phi \right. \\
& \left. -
\int_\om b(x)u_+^{r(x)-1}\phi\right\} +
\left\{ \int_\om|\grad u|^{p(x)}+
\int_{\Q} 
\frac{|u(x)-u(y)|^{p(x,y)}}{|x-y|^{N+sp(x,y)}} \right.\\
& \left. -\la
\int_\om a(x)u_+^{1-\ga(x)}-
\int_\om b(x)u_+^{r(x)} \right\} -
\int_{\om_2}|\grad u|^{p(x)}\\
    &-
\int_{\om_2}\ve|\na u|^{p(x)-2}\na u\na \phi+
\int_{\Q} 
\frac{W(x,y)}{|x-y|^{N+sp(x,y)}}\\
    &+\la
\int_{\om_2} a(x)u_+^{-\ga(x)}(u+\ve\phi)+
\int_{\om_2} b(x)u_+^{r(x)-1}(u+\ve\phi)\\
    \le& \ve I_1-
\int_{\om_2}\ve|\na u|^{p(x)-2}\na u\na \phi+
\int_{\Q} 
\frac{U(x,y)((u+\ve\phi)_-(x)-(u+\ve\phi)_-(y))}
{|x-y|^{N+sp(x,y)}}\\
    &+
\int_{\om_2} b(x)u_+^{r(x)-1}(u+\ve\phi),
\end{align*}
where 
\begin{equation*}
\begin{aligned}
I_1 =
\int_\om&|\grad u|^{p(x)-2}\na u\na\phi+
\int_{\Q} 
\frac{U(x,y)(\phi(x)-\phi(y))}{|x-y|^{N+sp(x,y)}} \\
&-\la
\int_\om a(x)u_+^{-\ga(x)}\phi -
\int_\om b(x)u_+^{r(x)-1}\phi.
\end{aligned}
\end{equation*}


\noindent Using above together with the estimates obtained for $W(x,y)$ gives 
    \begin{align*}
&0  \le \ve I_1-
\int_{\om_2}\ve|\na u|^{p(x)-2}\na u\na \phi+
\int_{\om_2} b(x)u_+^{r(x)-1}(u+\ve\phi) \\
&\quad -2
\int_{\om_2\times \om_1} 
\frac{U(x,y)(u+\ve\phi)(x)}{|x-y|^{N+sp(x,y)}} -2
\int_{\om_2\times \om^c} 
\frac{U(x,y)u(x)(u+\ve\phi)(x)}{|x-y|^{N+sp(x,y)}}\\
& \quad  -
\int_{\om_2\times \om_2} 
\frac{U(x,y)((u+\ve\phi)(x)-(u+\ve\phi)(y))}
{|x-y|^{N+sp(x,y)}}\\
    & \le  \ve I_1-
\int_{\om_2}\ve|\na u|^{p(x)-2}\na u\na \phi+
\int_{\om_2} b(x)u_+^{r(x)-1}(u+\ve\phi) \\
& \quad -2 \int_{\om_2\times \om_1} 
\frac{U(x,y)u(x)}{|x-y|^{N+sp(x,y)}} -2
\int_{\om_2\times \om^c} 
\frac{|u(x)|^{p(x,y)}}{|x-y|^{N+sp(x,y)}} \\
& \quad  - \int_{\om_2\times \om_2} 
\frac{|u(x)-u(y)|^{p(x,y)}}{|x-y|^{N+sp(x,y)}}-\ve
\left (2
\int_{\om_2\times \om_1} 
\frac{U(x,y)\phi(x)}{|x-y|^{N+sp(x,y)}}\right.\\
    & \quad \left. 2
\int_{\om_2\times \om^c} 
\frac{|u(x)|^{p(x,y)-2}u(x)\phi(x)}{|x-y|^{N+sp(x,y)}}+
\int_{\om_2\times \om_2} 
\frac{U(x,y)(\phi(x)-\phi(y))}{|x-y|^{N+sp(x,y)}}
\right )\\
    & \le  \ve I_1-\ve
\int_{\om_2}|\na u|^{p(x)-2}\na u\na \phi+
\int_{\om_2} b(x)u_+^{r(x)-1}(u+\ve\phi)\\
& \quad  -2
\int_{\om_2\times \om_1} 
\frac{U(x,y)u(x)}{|x-y|^{N+sp(x,y)}}
-2\ve
\left (
\int_{\om_2\times \om_1} 
\frac{U(x,y)\phi(x)}{|x-y|^{N+sp(x,y)}} \right. \\
& \quad \left. +
\int_{\om_2\times \om^c} 
\frac{|u(x)|^{p(x,y)-2}u(x)\phi(x)}
{|x-y|^{N+sp(x,y)}} + \frac{1}{2}
\int_{\om_2\times \om_2} 
\frac{U(x,y)(\phi(x)-\phi(y))}
{|x-y|^{N+sp(x,y)}}
\right )\\
    & \le   \ve I_1-\ve
\int_{\om_2}|\na u|^{p(x)-2}\na u\na \phi+\ve
\int_{\om_2} b(x)u_+^{r(x)-1}\phi \\
& +\ve\|b\|_{L^{
\frac{p^*(x)}{p^*(x)-r(x)}}(\om)}\ve^{r^--1}\Big (
\int_{\om_2}|\phi|^{p^*(x)}\Big )^{
\frac{r(x)}{p^*(x)}}\\
    &+2\ve\Big (
\int_{\om_2\times \om_1} 
\frac{|u(x)-u(y)|^{p(x,y)}}{|x-y|^{N+sp(x,y)}}\Big )^{
\frac{p(x,y)-1}{p(x,y)}}\Big (
\int_{\om_2\times \om_1} 
\frac{|\phi(x)|^{p(x,y)}}{|x-y|^{N+sp(x,y)}}\Big )^{
\frac{1}{p(x,y)}}\\
    &-2\ve\left[\Big (
\int_{\om_2\times \om^c} 
\frac{|u(x)|^{p(x,y)}}{|x-y|^{N+sp(x,y)}}\Big )^{
\frac{p(x,y)-1}{p(x,y)}}\Big (
\int_{\om_2\times \om^c} 
\frac{|\phi(x)|^{p(x,y)}}{|x-y|^{N+sp(x,y)}}\Big )^{
\frac{1}{p(x,y)}}\right.\\
    &+\Big (
\int_{\om_2\times \om_1} 
\frac{|u(x)-u(y)|^{p(x,y)}}{|x-y|^{N+sp(x,y)}}\Big )^{
\frac{p(x,y)-1}{p(x,y)}}\Big (
\int_{\om_2\times \om_1} 
\frac{|\phi(x)|^{p(x,y)}}{|x-y|^{N+sp(x,y)}}\Big )^{
\frac{1}{p(x,y)}}\\
    &\left.+
\frac{1}{2}\Big (
\int_{\om_2\times \om_2} 
\frac{|u(x)-u(y)|^{p(x,y)}}{|x-y|^{N+sp(x,y)}}\Big )^{
\frac{p(x,y)-1}{p(x,y)}}\Big (
\int_{\om_2\times \om_2} 
\frac{|\phi(x)-\phi(y)|^{p(x,y)}}
{|x-y|^{N+sp(x,y)}}\Big )^{
\frac{1}{p(x,y)}}\right].
    \end{align*}
    As the measure of the set 
    $\om_2:=\{x\in\om:u+\ve\phi\le0\}$
tends to $0$
as $\ve\to0.$
Therefore dividing by $\ve$
and letting $\ve\to0$
yields
\begin{align}\label{eq17}
\int_\om|\grad u|^{p(x)-2}\na u\na\phi+
&\iint_{\mathbb{R}^{2N}}
\frac{|u(x)-u(y)|^{p(x,y)-2}(u(x)-u(y))
(\phi(x)-\phi(y))}{|x-y|^{N+sp(x,y)}}\nonumber\\
    &-\la
\int_\om a(x)u_+^{-\ga(x)}\phi -
\int_\om b(x)u_+^{r(x)-1}\phi\ge0.
    \end{align}
Replacing $\phi$
by $-\phi$
in \eqref{eq17} gives the reverse inequality.
Hence, $u$
is a positive weak solution of \eqref{eq2} 
and therefore also a positive weak 
solution of \eqref{eq1}.


For $v\in \Nh^-,$
by applying Lemmas~\ref{lm9} and 
\ref{lm:min_neg}, 
together with similar arguments as above, 
we conclude that $v$
is a positive weak solution of \eqref{eq1}.
\end{proof}

\end{document}
